\documentclass[10pt,a4paper]{amsart}
\usepackage[margin=19mm]{geometry}
\usepackage{amsmath,amssymb,mathtools}
\usepackage[scr=rsfs,cal=euler]{mathalfa}
\usepackage{xfrac}
\usepackage[unicode,hidelinks,bookmarks=false]{hyperref}
\newcommand{\ie}{i.e.\ }
\newcommand{\cf}{cf.\ }
\newcommand{\resp}{respectively\ }
\newcommand{\Subsection}{\subsection}
\providecommand{\nobreakdash}{\nobreak\hbox{-}}
\setlength{\emergencystretch}{2em}
\multlinegap0pt
\usepackage[shortlabels]{enumitem}
\usepackage{amssymb}
\usepackage{mathtools}
\usepackage{tikz}
\usepackage{xcolor}
\newtheorem{lemma}{Lemma}
\newtheorem{proposition}{Proposition}
\newtheorem{corollary}{Corollary}
\usepackage{cleveref}
\crefname{lemma}{Lemma}{Lemmas}
\crefname{proposition}{Proposition}{Propositions}
\crefname{corollary}{Corollary}{Corollarys}
\newcommand{\Cern}{\operatorname{Cern}}
\newcommand{\M}{\mathcal M}
\newcommand{\eps}{\boldsymbol\varepsilon}
\title{The Class Edge-Reconstruction Number of a\\ Maximal Planar Graph Is One or Two}
\author{Sergey Ivanov}
\date{Source: arXiv:2609.02389v1 (CC BY 4.0)}
\begin{document}
\maketitle

\noindent\emph{Source and licence.} The Class Edge-Reconstruction Number of a\\ Maximal Planar Graph Is One or Two. Version arXiv:2609.02389v1 (CC BY 4.0), distributed by arXiv under the Creative Commons Attribution 4.0 International licence, http://creativecommons.org/licenses/by/4.0/, as recorded in the arXiv licence metadata for this exact version and shown on the abstract page https://arxiv.org/abs/2609.02389v1. Full licence notice: \texttt{licenses/arxiv-2609.02389-CC-BY-4.0.txt}.\par\smallskip

\noindent\textit{Editorial scope.} Counted unit: the article's proposition prop:cubic (source0.tex lines 366--369). The article's complete original proof of it is delivered with it (source0.tex lines 821--938, one \texttt{proof} environment of 118 lines, which the article places away from the statement). No mathematics is written, repaired or re-derived: the statement and the proof are the article's own lines, in the article's own notation, and no retained line is substituted or re-flowed.  Retained uncounted as the source-local context the proof needs: lines 266--270; lines 587--590; lines 603--607; lines 616--620; lines 686--706.  Nothing was omitted from any retained span, so the package carries no omission record.  No retained line contains a citation, so the wrapper carries no bibliography and no bibliography entry.  No retained line contains a figure environment, an image-inclusion command or drawing code, so no image is delivered and the package declares no asset dependency.  Editorial formatting only, in addition: an AMS \texttt{amsart} preamble that restates the article's own declarations and macros used by the retained lines, namely the environments \texttt{lemma}, \texttt{proposition}, \texttt{corollary} on separate counters, together with the standard packages those lines need.  The retained environments print in this document as Lemma 1, Proposition 1, Corollary 1 in order of appearance, whereas the article prints the same items with the numbering of its own full document.  The complete native source of the article as distributed in the arXiv e-print for this exact version is bundled unchanged as \texttt{sources/209194/source0.tex}.
\begin{lemma}[Two-parent lemma]\label{lem:two-parent}
If \(e\) is a flippable edge of a maximal planar graph \(G\), then \(G-e\)
is three-connected.  Its only maximal-planar parents are \(G\) and the flip
mate \(G^e\), up to isomorphism.
\end{lemma}

\begin{lemma}[Adjacent cubic vertices]\label{lem:adjacent-cubic}
If two degree-three vertices of a simple maximal planar graph are adjacent,
then the graph is \(K_4\).
\end{lemma}

\begin{corollary}\label{cor:cubic-independent}
In a maximal planar graph other than \(K_4\), the cubic vertices form an
independent set.  If \(v\) is cubic with neighbours \(a,b,c\), then
\(G-v\) is maximal planar and \(abc\) is one of its faces.
\end{corollary}

\begin{lemma}[Small orders]\label{lem:small-orders}
The maximal planar graphs on three, four, and five vertices are unique up to
isomorphism, namely \(K_3\), \(K_4\), and \(K_5-e\).  Each has class
edge-reconstruction number one.
\end{lemma}

\begin{lemma}[Residual identity]\label{lem:residual}
Suppose the same card has completions \(G=C+e\) and \(H=C+e'\), and set
\[
                  Q=A_G(e)-A_H(e').
\]
If \(G-g\cong H-h\), then
\[
                  A_G(g)-A_H(h)=Q.                 \tag{C.1}\label{eq:residual}
\]
Write \(Q=P-N\), where \(P,N\) have disjoint supports and nonnegative
integer entries.  Their common total mass \(r\) is zero, one, or two.  If
every edge card of \(G\) has a degree-sequence match among the edge cards of
\(H\), then:
\begin{enumerate}[label=\textup{(\roman*)},leftmargin=2.2em]
\item if \(r=2\), every edge of \(G\) has the same unordered endpoint-degree
pair;
\item if \(r=1\), the vertices of one fixed degree form a vertex cover of
\(G\);
\item if \(r=0\), then \(Q=0\).
\end{enumerate}
\end{lemma}

\begin{proposition}[The cubic case]\label{prop:cubic}
A maximal planar graph containing a vertex of degree three is identified by
at most two selected edge cards.
\end{proposition}

\begin{proof}[Proof of \cref{prop:cubic}]
The graphs of order at most five are covered by
\cref{lem:small-orders}, so assume \(n\ge6\).  Suppose for a contradiction
that \(\Cern_{\M}(G)>2\).  Let \(v\) be cubic, with cyclic neighbourhood
\(N(v)=\{a,b,c\}\).

For each \(f\in\{ab,bc,ca\}\), let \(z_f\) be the third vertex of the face
on the side of \(f\) opposite \(v\).  The vertex \(z_f\) is outside
\(\{v,a,b,c\}\): otherwise those four vertices form \(K_4\), whose faces
and the degree of \(v\) leave no room for the remaining vertices.  Since
\(v\) has only the neighbours \(a,b,c\), the edge \(vz_f\) is absent.
Thus every \(f\) is flippable.  Put
\[
             C_f=G-f,\qquad H_f=G-f+vz_f.
\]
By \cref{lem:two-parent}, \(G\) and \(H_f\) are the only parents of
\(C_f\).  If \(H_f\cong G\), then \(C_f\) already identifies \(G\),
contrary to our assumption.  Hence \(H_f\not\cong G\).

Fix a boundary edge \(f\).  For every edge \(g\ne f\), the selected pair
\(C_f,G-g\) does not identify \(G\).  Its only possible nonisomorphic parent
after the first card is \(H_f\), so \(H_f\) contains a card isomorphic to
\(G-g\); if \(G-g\cong C_f\), the same conclusion holds with the required
multiplicity.  For \(g=f\), the common card itself is
\(G-f=H_f-vz_f\).  Consequently, for every edge \(g\) of \(G\), there is an
edge \(h\) of \(H_f\) such that
\[
                         G-g\cong H_f-h.           \tag{E.1}\label{eq:allmatch}
\]

For \(f=xy\), write \(s=d_G(z_f)\).  Since deleting \(xy\) lowers its two
endpoint degrees, while the competing completion inserts \(vz_f\), the
residual in \cref{lem:residual} is
\[
 Q_f=\eps_{d(x)}+\eps_{d(y)}-\eps_4-\eps_{s+1}.    \tag{E.2}\label{eq:cubicQ}
\]
After cancelling equal positive and negative units, let \(r_f\) be the mass
of the positive part.

We first show \(r_f\ne2\).  By \eqref{eq:allmatch} and
\cref{lem:residual}, mass two would force every edge of \(G\) to have one
fixed unordered endpoint-degree pair.  If the two degrees differ, every edge
joins the two corresponding degree classes and \(G\) is bipartite,
contradicting its triangular faces.  If they agree, connectedness makes
\(G\) regular; because \(G\) has a cubic vertex, Euler's formula then gives
\(G=K_4\), contrary to \(n\ge6\).  Hence \(r_f\in\{0,1\}\) for all three
boundary edges.

Suppose first that \(r_f=0\) for all \(f\in\{ab,bc,ca\}\).  Equation
\eqref{eq:cubicQ} says that each of the pairs
\[
           \{d(a),d(b)\},\quad \{d(b),d(c)\},\quad \{d(c),d(a)\}
\]
contains \(4\).  At least two of \(a,b,c\) therefore have degree four.
In the maximal planar graph \(T=G-v\), supplied by
\cref{cor:cubic-independent}, those two vertices are adjacent cubic
vertices.  \Cref{lem:adjacent-cubic} gives \(T=K_4\), hence \(n=5\), a
contradiction.

We may choose a boundary edge \(f\) with \(r_f=1\), and write its reduced
residual as
\[
                         Q_f=\eps_p-\eps_q.          \tag{E.3}\label{eq:massone}
\]
By \eqref{eq:allmatch} and \cref{lem:residual}, the degree-\(p\) vertices
meet every edge of \(G\).  We cannot have \(p=3\): by
\cref{cor:cubic-independent}, the cubic vertices would then be an independent
vertex cover, making both it and its complement independent and making
\(G\) bipartite.  Since \(v\) is not a degree-\(p\) vertex, each of the
edges \(va,vb,vc\) forces
\[
                         d(a)=d(b)=d(c)=p.           \tag{E.4}\label{eq:pneighbors}
\]
If \(p=4\), then \(G-v\) again has adjacent cubic vertices and is \(K_4\),
contrary to \(n\ge6\).  Thus \(p\ge5\).

Restart the construction for each boundary edge \(xy\), using its own flip
mate \(H_{xy}\), and apply \eqref{eq:cubicQ} separately.  By
\eqref{eq:pneighbors}, its uncancelled positive part is \(2\eps_p\), while
the negative part is \(\eps_4+\eps_{d(z_{xy})+1}\).  Since \(p\ne4\) and
residual mass two is impossible, one negative unit must cancel a copy of
\(\eps_p\).  Therefore
\[
       d(z_{xy})=p-1,
       \qquad Q_{xy}=\eps_p-\eps_4                  \tag{E.5}\label{eq:allQ}
\]
for each of \(ab,bc,ca\).

Let \(p=5\), and inspect the cyclic link of \(a\).  It contains the
consecutive segment \(b,v,c\).  On the complementary \(c\)-to-\(b\) arc,
the first and last vertices are \(z_{ac}\) and \(z_{ab}\), both of degree
four by \eqref{eq:allQ}.  The degree-five vertices cover every edge, so two
non-degree-five vertices cannot be adjacent.  If
\(z_{ac}=z_{ab}\), then the link has length four; if they are distinct,
they require at least one intermediate vertex and the link has length at
least six.  Both alternatives contradict \(d(a)=5\).

It remains that \(p\ge6\).  Fix \(f=ab\), write \(z=z_{ab}\), and put
\(H=H_f\).  The card \(G-va\) must occur in \(H\), so
\cref{lem:residual} and \eqref{eq:allQ} require an edge \(h\in E(H)\) with
\[
  A_H(h)=A_G(va)-Q_f
        =(\eps_3+\eps_p)-(\eps_p-\eps_4)
        =\eps_3+\eps_4.                              \tag{E.6}\label{eq:34edge}
\]
Thus \(H\) would contain a degree-three--degree-four edge.  It does not.
Indeed, the degree-\(p\) vertices cover every edge of \(G\), while the flip
changes only these degrees:
\[
  a,b:p\mapsto p-1,\qquad v:3\mapsto4,
  \qquad z:p-1\mapsto p.
\]
Every degree-three vertex of \(H\) is therefore an unchanged cubic vertex of
\(G\), and all its neighbours in \(H\) have degree at least
\(p-1\ge5\).  The new degree-four vertex \(v\) has neighbours of degrees
\(p-1,p-1,p,p\).  No degree-three--degree-four edge exists, contradicting
\eqref{eq:34edge}.  This final contradiction proves the proposition.
\end{proof}

\end{document}
